\section{Introduction}
\label{sec:intro}

Industrial predictive maintenance increasingly requires models able to \emph{learn continually} on the
equipment itself, without a round-trip to a server: operating conditions drift (wear, regime or machine
change), and a frozen model degrades. \gls{CL} answers this
need~\cite{DeLange2021Survey,Hurtado2023CLPdM}, yet its deployment on an \gls{MCU} remains
constrained: \gls{SRAM} in the range of a hundred kilobytes, no neural accelerator, a strict latency budget.

Against these constraints, INT8 quantization is the best-established reflex of the embedded
literature~\cite{Ravaglia2021QLRCL,Capogrosso2023TinyML}. It is almost always justified by a threefold
argument --- less memory, less time, less energy --- that is seldom verified \emph{on the target} for each
of the three terms separately. This is precisely what we measure here, and the three terms do not behave
alike.

This work targets a \textbf{triple gap} that few studies close simultaneously: (Gap~1) validation on real
industrial time series; (Gap~2) operation under a \emph{strict and measured} memory and latency budget ---
here \SI{256}{\kilo\byte} of SRAM and \SI{100}{\milli\second} per inference + update cycle; (Gap~3) INT8
quantization compatible with incremental training. We instantiate these three axes on an \gls{EWC}
classifier~\cite{Kirkpatrick2017EWC} ported to C on a NUCLEO-F439ZI board (Cortex-M4
\SI{180}{\mega\hertz}, \SI{256}{\kilo\byte} SRAM, no \gls{NPU}), evaluated on two industrial datasets with
contrasting continual scenarios.

\paragraph{Contributions.} (i)~A \emph{paired} PC$\leftrightarrow$board protocol (same seed, same sample
order, same normalization) establishing an \emph{exact} FP32 parity in the frozen protocol (\num{1.000}),
near-exact online (from \num{0.9626} to \num{0.9887} across the three cells), for an inference + learning
latency
$\ll$~\SI{100}{\milli\second}. (ii)~The demonstration, \emph{on real hardware}, of the collapse of naive
\gls{PTQ} of the binary EWC head: F1 falls from $\approx\!\num{0.92}$ to \num{0.138} and \num{0.1337} even
though weight \gls{RAM} is indeed divided by~4. (iii)~The isolation of the cause through a \emph{bit-exact
emulated} ablation, then the \emph{recovery measured on-board} with a calibrated-scale kernel: F1 returns
to \num{0.9173} and \num{0.8995}, with a frozen parity of \num{1.000} against the emulator and no \gls{CRC}
error. (iv)~Two complementary quantization axes: the \emph{moment}, where \gls{QAT} training adds no decisive
gain over a properly calibrated PTQ, and the \emph{depth}, where the boundary depends on the metric chosen
--- ternary in AUROC, INT4 in on-board measured F1 --- and where the memory gain exists only with genuine
bit-packing. (v)~A complete system budget measured on the same
board, which \emph{decouples the three promises of INT8}: weight RAM is indeed divided by~4, but total
system RAM is unchanged, latency increases by about half, and energy does not move.

This last point is the guiding thread of the article. On this target, \textbf{INT8 is justified by memory,
neither by speed nor by energy} --- and even then, by a memory whose nature must be stated. Throughout the
text and the captions we systematically distinguish what is \emph{measured on-board} from what is
\emph{bit-exactly emulated on PC}, and we explicitly leave "to be measured" any cell not yet streamed,
rather than estimating it.

\paragraph{Organization.} Section~\ref{sec:related} positions the work, Section~\ref{sec:method} describes
the model, the paired protocol and the quantization schemes, Section~\ref{sec:setup} the setup. Results
come in three stages: parity, collapse and recovery (Section~\ref{sec:results}), moment and depth
(Section~\ref{sec:quant-axes}), system budget (Section~\ref{sec:system}). Section~\ref{sec:discussion}
discusses both paradoxes and the limitations.
